\documentclass{article}
\usepackage[utf8]{inputenc}
\usepackage[T1]{fontenc}
\usepackage[margin=2cm]{geometry}
\usepackage{hyperref}
\sloppy
\begin{document}
\section*{Historical source}
Encyclopaedia Britannica scan, source row 1214. Parse 3.5 extracted the text. Super 3.5 VL corrected it from an enlarged crop of each block. It can still contain recognition errors.
Page headers, footers, and page numbers are removed.

\subsection*{Untitled continuation (E0095)}
\textit{Continues an entry from an earlier page.}

DAM

Damps.

thickened the air in the pit, and then it would appear again; they lighted their candles at it sometimes when they went out; and so in this pit it did no further harm.

In another pit, however, it soon appeared, and at last produced a most terrible explosion. This was occasioned by one of the workmen going imprudently down with a lighted candle, after a cessation of work for some days, and the force exerted by it seemed equal to that of gun-powder.

The formation of unintractable air in mines cannot be attributed to any variation of atmospheric or there is no natural process with which we are acquainted, by which such a change could be accomplished. In one instance, however, we have an example of a fire-damp being produced, not only without any considerable flagration of atmospheric air, but where there is the belt circulation imaginable. This is in large bellows used in metallurgic works, which are sometimes burst by an explosion of inflammable matter proceeding from the rancid matters with which the leather is greased. Dr Priestley has shown, that inflammable air is composed of pure elementary fire, charcoal or phlogiston, and a little water; and that this composition may take place even in vacuo. All the materials are to be met with in the bowels of the earth. Coal, a bituminous substance, is abundantly able to supply the phlogiston; the natural moisture of the earth affords water, and the heat, however gentle, which constantly exists in the bowels of the earth, may be sufficient to produce a quantity of inflammable air, which gradually accumulating in those places where there is not a constant stream of atmospheric air to carry it off, will soon produce the dreadful effects already mentioned.

A much more important consideration than the formation of damps, however, is the proper method of avoiding their pernicious effects. The inflammability of one kind affords an easy method of preventing it from accumulating, viz. by letting fire to it. This may be done with safety, unless it has been suffered to go too far before the experiment is made: for the inflammable air, being much lighter than any other kind, will naturally rise to the top; so that a man, lying flat on the ground to avoid the force of the explosion, and holding up a lighted candle fixed upon a pole, may at once free the mine from such a troublesome guest. But where it has been allowed to accumulate in too great quantity, so that this method cannot be used, or in the other kind, which is not inflammable, the method commonly practised is to produce a continual circulation of air as much as possible through all parts of the mine. To procure this, they make a perpendicular opening, which they call a flank or flight, so that the mine may have two or more openings; and thus by reason of the difference of temperature between the open atmosphere and that in the mine, there is a continual draught of air through them both. This current will always be stronger in proportion to the difference between the external atmosphere and that of the mine; and likewise in proportion to the difference between the depth of the two shafts. But as the temperature of the atmosphere is variable, it happens, at certain seasons of the year, that there is not a sufficient difference between that

of the atmosphere and in the mine to produce the necessary circulation. This happens principally in the spring and autumn; at which seasons it is necessary to light fires in the shafts, which are always efficacious for the purpose desired.

Among the other uses to which dephlogisticated air might be applied, Mr Cavallo reckons that of rescuing people from the dangerous effects of damps in mines, and other subterraneous places. "If a large bladder," says he, "into which a solution of lime in water is introduced, be filled with dephlogisticated air, and a small wooden or glass pipe be adapted to its neck, a man may hold that pipe in his mouth, and may breathe the dephlogisticated air; and thus equipped he may enter into the subterraneous places, amidst the various noxious fluids contained in them. A large bladder of dephlogisticated air will serve for above a quarter of an hour, which is a length of time sufficient for various purposes; besides, if longer time is required to be spent in these places, a person may have two or more bladders of dephlogisticated air along with him, and may shift as soon as the air of one is contaminated. Without the necessity of any more complicated apparatus, the bladders full of dephlogisticated air may be kept stopped by putting corks into the glass or wooden pipes that are tied to their necks. This air might also be used for diving-bells.

\subsection*{DAMSEL (E0096)}
DAMSEL, from the French damoisel or damoisel, an appellation anciently given to all young people of either sex, that were of noble or gentle extraction, as the sons and daughters of princes, knights, and barons: thus we read of Damel Pepin, Damel Louis le Gros, Damel Richard prince of Wales.

From the sons of kings this appellation first passed to those of great lords and barons, and at length to those of gentlemen who were not yet knights.

At present damsel is applied to all maids or girls not yet married, provided they be not of the vulgar.

\subsection*{DAN (E0097)}
DAN, or JOR-DAN, which last literally denotes “the river Dan;” to named from the people where it has its source, which is a lake called Phiala, from its round figure, to the north of its apparent rising from the mountain Panium or Paneum, as was discovered by Philip, Tetrarch of Trachonites; for on throwing light bodies into the Phiala, he found them to emerge again at Paneum (Josephus). From Paneum it runs in a direct course to a lake called Samachonitis; as far as which it is called Jordan the Lesser; and thence to the lake Gennesareth, or of Tiberias, where it comes increased by the lake Samachonitis and its springs, and is called the Greater Jordan; continuing its direct course southwards, till it falls into the Asphaltites.

DAN (anc. geog.), a town to the west of the source of the Jordan; formerly called Laish (Joshua, Judges, Josephus). This was the north, as Beersheba was the south, boundary of the Israelites; as appears from the common expression in Scripture, from Dan to Beersheba. At Dan Jeroboam erected one of the golden calves (1 Kings xii.).

calves (1 Kings xii. 29).
DAN, the tribe, extended itself westward of Judah,
and was terminated by Azotus and Dora on the Medi-
terranean (Joppa).

\subsection*{DANAE (E0098)}
DANAE, in antiquity, a coin somewhat more than an obolus, used to be put into the mouths of the dead, to pay their passage over the river Acheron.


\section*{Source record}
\begin{description}
\item[Scan ID] \texttt{britannica:309f765b1ecc83e0a7f7762e44b4a927a7555443:train:1214}
\item[Image file] \texttt{page-1214-17d180a99e0a.jpg}
\item[Image SHA-256] \texttt{17d180a99e0a48546a555d98cd7fdf15bb8f8ea4b751c50e02536eba79c2473a}
\item[Dataset] NationalLibraryOfScotland/encyclopaedia\_britannica\_illustrated at revision \texttt{309f765b1ecc83e0a7f7762e44b4a927a7555443} (CC0-1.0)
\item[Extraction] model \texttt{nvidia/nemotron-3.5-parse}, raw reply \texttt{extract/raw/page-1214-17d180a99e0a.json}
\item[Raw reply SHA-256] \texttt{853588341eb038a7e25c4f31170c11ef5726149bf97f0e044d9dcf3f6d236d9b}
\item[Correction] model \texttt{nvidia/Nemotron-3.5-Super-VL}, block crops, raw replies \texttt{correct/raw/crop\_nothink/page-1214-17d180a99e0a-b*.json}; parse: 1, super\_vl: 15
\item[Running header] DAM [ 658]
\end{description}
\end{document}
